\section{Component Interactions}
\label{sec:component_interactions}

The system components interact synchronously across well-defined boundaries. Table \ref{tab:component_responsibility} summarizes the language and core responsibility of each component in the pipeline.

\begin{table}[H]
\centering
\caption{Component Responsibility Summary}
\label{tab:component_responsibility}
\small
\begin{tabularx}{\textwidth}{>{\hsize=0.6\hsize}X >{\hsize=0.6\hsize}X >{\hsize=1.8\hsize}X}
\toprule
\textbf{Component} & \textbf{Language} & \textbf{Core Responsibility} \\
\midrule
Workload Gen. & Python & Synthesize controlled traffic (Poisson arrivals) and simulate varying network load. \\
\addlinespace
Sequencer & Rust & Validate, order, and batch L2 transactions using policies like FCFS; enforce timeout triggers. \\
\addlinespace
Executor & Rust & Receive gRPC batches, execute a transfer-centric state transition over a RocksDB Sparse Merkle Tree, and invoke the Prover. \\
\addlinespace
Prover & Rust & A subsystem of the Executor using \texttt{risc0\_prover} (or mock delays) to generate cryptographic validity artifacts. \\
\addlinespace
Submitter & Rust & Handle L1 Ethereum Adapter interactions, manage DA formatting (Calldata/Blobs), and ensure reliable transaction broadcast via a Saga loop. \\
\addlinespace
Contracts & Solidity & L1 bridge; validates DA commitments, verifies proofs via \texttt{MockVerifier} or Groth16, and finalizes the batch. \\
\bottomrule
\end{tabularx}
\end{table}

The interactions ensure a strict cryptographic link. The Sequencer's transaction handling is isolated from batch creation. When a batch is sealed, it is sent to the Executor via gRPC. The Executor maintains a trace of the execution and invokes the prover backend. The resulting enriched payload (DA commitment and proof) is broadcast via a stream to the Submitter, which handles the final L1 gas estimation and submission.
